\section{The joint impact record and physical sampling cuts}
\label{sec:records}
The endpoint parametrization gives more than a count or a marginal roof.
It reconstructs all positions, directions and collision times in the
window. We describe its onset law explicitly.

Let
\begin{equation}\label{eq:paraboloid}
 \mathcal U=\{(z,\eta)\in\R^2\times\R:
                   |z|<1,\ 0<\eta<1-|z|^2\},\qquad
 \dd\mathsf m=\frac2\pi\one_{\mathcal U}\dd z\dd\eta.
\end{equation}
This is a probability measure. For $z\in\R^2$, set
$(U,V)^t=H_{j,R}^{-1/2}z$ and define
\begin{align}\label{eq:limiting-bridge}
 Y_i(z)&=\frac{\sinh((j-i)\chi_R)}{\sinh(j\chi_R)}U
       +\frac{\sinh(i\chi_R)}{\sinh(j\chi_R)}V,\qquad 0\le i\le j,\\
 Q_i(z)&=\frac{Y_i^2+Y_{i+1}^2}{R}
                    +\frac{(Y_{i+1}-Y_i)^2}{g_R},\qquad 0\le i<j.\notag
\end{align}
These quadratic functions satisfy $Q_i\ge0$ and
$\sum_{i=0}^{j-1}Q_i=|z|^2$.

\begin{theorem}[Marked onset record]\label{thm:record}
Condition on $N_{jg_R+\varepsilon}=j+1$. The orientation $e$ is exactly
uniform on the six nearest lattice vectors. In that orientation, use
the exact Morse coordinates $w$ of the endpoints and the first-impact
time $r$, and set $z=w/\sqrt{2\varepsilon}$, $\eta=r/\varepsilon$.
Their law has total variation distance at most $C\varepsilon$ from
$\mathsf m$, uniformly in $j\ge1$ and $R\in K$.

On these same coordinates, the transverse impact positions and impact
times $t_i$ satisfy
\begin{equation}\label{eq:record-asymptotic}
 \frac{y_i}{\sqrt{2\varepsilon}}=Y_i(z)+O(\varepsilon),\qquad
 \frac{t_i-ig_R}{\varepsilon}
       =\eta+\sum_{h=0}^{i-1}Q_h(z)+O(\varepsilon),
 \qquad 0\le i\le j,
\end{equation}
with errors uniform in the index, itinerary length, radius and
$(z,\eta)\in\mathcal U$. Both velocity traces are obtained from the
same bridge by the chord and reflection formulas.

Consider any fixed finite number of samples at
$t^{\rm s}=kg_R+\kappa\varepsilon\in[0,jg_R+\varepsilon]$, with
$0\le k\le j$ and bounded fixed $\kappa$. In the limiting record, the
sample precedes impact $k$ exactly on
\begin{equation}\label{eq:sample-cut}
 \eta+\sum_{h<k}Q_h(z)>\kappa.
\end{equation}
The joint probabilities of the resulting before/after cuts converge
with error $O(\varepsilon)$, uniformly in $j$ and $R$. Thus the
sampling-time interfaces are retained, not suppressed by pointwise
differentiation away from event times.
\end{theorem}
\begin{proof}
Keep the $r$ integral in \eqref{eq:record-flux}. The changes
$w=\sqrt{2\varepsilon}z$, $r=\varepsilon\eta$ give exactly the domain
$\mathcal U$ and, up to its normalizing constant, density
$a_{j,R}(\sqrt{2\varepsilon}z)$. By \eqref{eq:morse-flux}, this density
is $1+O(\varepsilon)$ uniformly. Since
$\operatorname{vol}\mathcal U=\pi/2$, normalization proves the total
variation assertion. Rotation through multiples of $\pi/3$ preserves
phase volume and the event, proving the exact orientation weights.

The odd Morse inverse has expansion
$\Phi_{j,R}(w)=H_{j,R}^{-1/2}w+O(|w|^3)$. Insert it into
\eqref{eq:bridge-estimates}. This proves the first formula in
\eqref{eq:record-asymptotic}, with the stronger error
$O(\varepsilon(\rho^i+\rho^{j-i}))$. Expanding each chord with
\eqref{eq:one-quadratic} gives
$\ell_R(y_i,y_{i+1})=g_R+\varepsilon Q_i(z)
 +O(\varepsilon^2(\rho^i+\rho^{j-i})^2)$.
The sum of these error bounds is uniform, proving the time formula.
The identity for $\sum Q_i$ is the quadratic endpoint action evaluated
at $H_{j,R}^{-1/2}z$.

For clarity, in the pair's fixed global transverse direction the leading
outgoing transverse velocity at impacts $i<j$, divided by
$\sqrt{2\varepsilon}$, is $(Y_{i+1}-Y_i)/g_R$.
At the final impact it is
$(Y_j-Y_{j-1})/g_R+2Y_j/R$; the initial incoming value is
$(Y_1-Y_0)/g_R-2Y_0/R$.
The longitudinal components and the unscaled velocities are given
exactly by the unit chord and specular formulas. Uniform analyticity
on the fixed arcs gives their corresponding Taylor bounds. Together
with the initial and terminal pieces in Proposition~\ref{prop:marked-flux},
this reconstructs the entire physical path.

Consecutive impacts are separated by at least $g_R$, so a sample within
$O(\varepsilon)$ of $kg_R$ can cross only the $k$th impact for small
$\varepsilon$. Formula \eqref{eq:record-asymptotic} gives
\eqref{eq:sample-cut}. The cut has derivative one in $\eta$.
For each fixed $z$, the set where its left side lies within $\delta$
of $\kappa$ has $\eta$-length at most $2\delta$. Its
$\mathsf m$-mass is therefore at most $4\delta$. Taking
$\delta=C\varepsilon$ and adding the $O(\varepsilon)$ density error
bounds the mismatch probability. A union bound proves the result for
any fixed finite set of cuts. Bounded Lipschitz tests of the scaled
positions and times can be added using their uniform pathwise errors.
No differentiability of a discontinuous sampling decoder is assumed.
\end{proof}

For the source-weighted conditional law, the orientation weights tend
to $e^{qw_{j,e}}/\sum_{e'}e^{qw_{j,e'}}$. Given an orientation, the
limiting law is still \eqref{eq:paraboloid}; its total variation error
is $O(\sqrt\varepsilon)$ uniformly for bounded sources, by
Lemma~\ref{lem:marks}. The integrated source normalization has the
stronger $O(\varepsilon)$ relative error of Theorem~\ref{thm:marked}
because its odd part cancels. The distinction between these two errors
is important for a full marked law.
